Une exploratrice emporte une horloge à pendule simple dont la tige (de masse négligeable) mesure \qty{\v{L}}{m}. Elle atterrit sur une planète où le champ gravitationnel est de \qty{\v{gp}}{m/s^2}. On suppose que le pendule oscille avec de petites amplitudes.

\begin{enumerate}[label=\alph*)]
    \item Quelle est la période du pendule sur Terre ($g = \qty{9.81}{m/s^2}$) ?
    \item Quelle est la période du pendule sur la planète ?
    \item Quelle longueur faudrait-il donner au pendule pour que l'horloge batte la même période sur la planète que sur Terre ? Faut-il allonger ou raccourcir la tige ?
\end{enumerate}
